\chapter{Contratos con el Cliente}\label{cap:contratos}

Un contrato es una versión PDF que firman primero el Cliente y después una
persona del despacho, cada uno con su certificado. El recorrido completo es:

\begin{enumerate}
  \item El despacho crea el contrato y solicita las firmas
        (\cref{sec:crear-contrato}).
  \item El Cliente revisa y firma desde \textbf{Mis contratos} (\cref{sec:cliente-firmar}).
  \item El despacho comprueba la firma del Cliente y firma desde
        \textbf{Por firmar} (\cref{sec:firmar-despacho}).
\end{enumerate}

\importante{Una versión en contrato ya no admite sello ni firma personal. Si
necesita cambiar el documento, agregue una versión nueva y solicite firmas
sobre ella.}

% ---------------------------------------------------------------------
\section{Antes de solicitar firmas}\label{sec:contrato-requisitos}

Verifique que se cumplen estas condiciones; si falta alguna, Qadra
muestra la lista \textbf{Antes de solicitar las firmas} y no crea la
solicitud.

\begin{table}[H]
  \centering
  \caption{Condiciones para solicitar firmas.}\label{tab:contrato-requisitos}
  \begin{tabularx}{\textwidth}{L{3.6cm} Y}
    \toprule
    \textbf{Quién o qué} & \textbf{Condición} \\
    \midrule
    Documento & Versión actual, en PDF y sin sellar. \\
    Cliente   & Cuenta creada por invitación, activa y con su firma vinculada
                en \textbf{Mi firma digital}. \\
    Firmante del despacho & Administrador o Litigante con firma vinculada y
                vigente, y asignado al expediente. \\
    Hoja de huella & Ambas partes firmaron a mano los ejemplares de la hoja de
                huella; se necesita su folio. \\
    \bottomrule
  \end{tabularx}
\end{table}

% ---------------------------------------------------------------------
\section{Crear un contrato y solicitar firmas}\label{sec:crear-contrato}

\paraque Convierte una versión PDF en contrato y la envía a firma.

\disponible Administrador y Litigante.

\pasos
\begin{enumerate}
  \item Abra la ficha del documento PDF (\cref{sec:ficha-documental}) y
        haga clic en \textbf{Crear contrato}.
  \item Revise el \textbf{Título del contrato} y haga clic en
        \textbf{Crear contrato}. Aparece \enquote{Contrato creado en
        borrador.}
  \item En \textbf{Solicitar firmas}, escriba el
        \textbf{Correo de la cuenta del Cliente}.
  \item En \textbf{Firma por el despacho}, elija la cuenta que firmará.
  \item Indique la \textbf{Vigencia en días} (de 1 a 30) y el
        \textbf{Folio de la hoja de huella firmada}.
  \item Haga clic en \textbf{Solicitar firmas}.
\end{enumerate}

\begin{figure}[H]
  \centering
  \capturapendiente{capturas/A-36-solicitar-firmas.png}%
    {Formulario Solicitar firmas con correo del Cliente, firmante, vigencia y folio · Litigante}
  \caption{Solicitud de firmas de un contrato.}
  \label{fig:manual-solicitar-firmas}
\end{figure}

\resultado Aparece \enquote{Firmas solicitadas. Consulta su avance en
Contratos del expediente.} El Cliente ya puede firmar.

% ---------------------------------------------------------------------
\section{Seguir un contrato}\label{sec:seguir-contrato}

\paraque Muestra el estado de los contratos del expediente y permite
descargar sus paquetes.

\disponible Administrador y Litigante; el Asistente legal solo consulta.

\pasos
\begin{enumerate}
  \item En el expediente, abra la pestaña \textbf{Contratos}.
  \item Haga clic en \textbf{Abrir} junto al contrato.
  \item Revise las \textbf{Solicitudes de firma} y la
        \textbf{Línea de tiempo}.
  \item Descargue lo que necesite: \textbf{Descargar acuse parcial} (cuando
        firmó el Cliente) o \textbf{Descargar contrato firmado} (cuando
        firmaron ambos).
\end{enumerate}

\begin{figure}[H]
  \centering
  \capturapendiente{capturas/A-37-detalle-contrato.png}%
    {Detalle de un contrato con su solicitud, línea de tiempo y botones de descarga · Litigante}
  \caption{Detalle de un contrato.}
  \label{fig:manual-detalle-contrato}
\end{figure}

\begin{table}[H]
  \centering
  \caption{Estados de una solicitud de firma.}\label{tab:estados-solicitud}
  \begin{tabularx}{\textwidth}{L{5cm} Y}
    \toprule
    \textbf{Estado} & \textbf{Qué significa} \\
    \midrule
    \textbf{Espera la firma del Cliente}  & El Cliente aún no firma. \\
    \textbf{Espera la firma del despacho} & El Cliente firmó; falta el despacho. \\
    \textbf{Formalizado}                  & Firmaron ambas partes. \\
    \textbf{Rechazado por el Cliente}     & El Cliente rechazó el contrato. \\
    \textbf{Declinado por el despacho}    & El despacho no firmó; la firma del Cliente se conserva. \\
    \textbf{Cancelado} / \textbf{Vencido} & La solicitud se cerró sin formalizar. \\
    \bottomrule
  \end{tabularx}
\end{table}

\begin{itemize}
  \item Para cancelar una solicitud abierta, haga clic en
        \textbf{Cancelar solicitud}, escriba el \textbf{Motivo} y haga clic
        en \textbf{Confirmar cancelación}.
  \item Un contrato en \textbf{Borrador} sin solicitud puede abandonarse con
        \textbf{Abandonar contrato}.
  \item Solo puede haber una solicitud abierta por contrato.
\end{itemize}

% ---------------------------------------------------------------------
\section{Firmar por el despacho}\label{sec:firmar-despacho}

\paraque Contrafirma el contrato después del Cliente y lo formaliza.

\disponible La cuenta elegida en \textbf{Firma por el despacho}
(Administrador o Litigante).

Tenga a la mano su clave privada cifrada (archivo \texttt{.pem}) y su
contraseña.

\pasos
\begin{enumerate}
  \item En el menú, haga clic en \textbf{Por firmar} y elija la solicitud.
  \item En \textbf{1. Revisa}, haga clic en
        \textbf{Cargar y comprobar el PDF} y lea el contrato.
  \item En \textbf{Comprueba la firma del Cliente}, haga clic en
        \textbf{Descargar acuse parcial}. Seleccione ese archivo en el
        validador y haga clic en \textbf{Validar paquete}.
  \item En \textbf{2. Confirma}, marque \textbf{Confirmo firmar por el
        despacho esta versión exacta, después de revisar la firma del
        Cliente.}
  \item En \textbf{3. Firma}, haga clic en \textbf{Preparar la firma}.
        Seleccione su \textbf{Clave privada cifrada}, escriba la
        \textbf{Contraseña de la clave} y haga clic en
        \textbf{Firmar localmente}.
  \item Haga clic en \textbf{Enviar firma}.
\end{enumerate}

\begin{figure}[H]
  \centering
  \capturapendiente{capturas/A-38-por-firmar.png}%
    {Por firmar con los pasos 1. Revisa, 2. Confirma y 3. Firma y el acuse parcial validado · Litigante}
  \caption{Firma del contrato por el despacho.}
  \label{fig:manual-por-firmar}
\end{figure}

\resultado Aparece \enquote{Firma del despacho registrada; el contrato quedó
formalizado.} Descargue el paquete con \textbf{Descargar contrato firmado}.

\begin{itemize}
  \item El botón de firma se habilita solo cuando el acuse parcial se valida
        sin fallas.
  \item Si intenta firmar antes que el Cliente, aparece \enquote{No es tu
        turno de firmar. Primero firma el Cliente y después el despacho.}
  \item Si no desea firmar, haga clic en \textbf{Declinar firma}.
  \item Si se pierde la respuesta, haga clic en
        \textbf{Consultar recibo de firma}; consultar no vuelve a enviar la
        firma.
\end{itemize}
